% ============================================================
% apress-style.tex — Stile Apress / Springer per libri tecnici
% ============================================================
% Caratteristiche: layout accademico pulito, font classici,
% capitoli con linea nera, stile formale e compatto.
% ============================================================

% === LINGUA ===
\usepackage[\BabelLang]{babel}
\usepackage[autostyle]{csquotes}

% === BIBLIOGRAFIA ===
\usepackage[
  backend=biber,
  style=numeric,
  sorting=nyt,
  urldate=long,
  maxbibnames=5,
  backref=true,
  backrefstyle=none
]{biblatex}

% === FONT (Apress: Computer Modern / Latin Modern classico) ===
\usepackage{fontspec}
\usepackage{amsmath}
\usepackage{amssymb}
\setmainfont{texgyretermes}[
  Extension=.otf,
  UprightFont=*-regular,
  BoldFont=*-bold,
  ItalicFont=*-italic,
  BoldItalicFont=*-bolditalic
]
\setsansfont{texgyreheros}[
  Extension=.otf,
  Scale=0.92,
  UprightFont=*-regular,
  BoldFont=*-bold,
  ItalicFont=*-italic,
  BoldItalicFont=*-bolditalic
]
\setmonofont{InconsolataN}[
  Extension=.otf,
  Scale=0.85,
  UprightFont=*-Regular,
  BoldFont=*-Bold
]

% === LAYOUT PAGINA (6.14 × 9.21 poll. — formato Apress standard) ===
\usepackage[
  paperwidth=6.14in,
  paperheight=9.21in,
  top=0.85in, bottom=0.85in,
  inner=0.95in, outer=0.7in,
  headheight=14pt,
  footskip=0.4in
]{geometry}

% === SPAZIATURA PARAGRAFI (Apress: indent tradizionale) ===
\setlength{\parindent}{1.5em}
\setlength{\parskip}{0pt plus 1pt}

% === COLORI (palette Apress — giallo/nero) ===
\usepackage[dvipsnames,svgnames,table]{xcolor}
\definecolor{apress-yellow}{HTML}{FFD700}
\definecolor{apress-dark}{HTML}{1A1A1A}
\definecolor{apress-gray}{HTML}{4A4A4A}
\definecolor{apress-lightgray}{HTML}{999999}
\definecolor{code-bg}{HTML}{F5F5F5}
\definecolor{code-border}{HTML}{CCCCCC}
\definecolor{tip-green}{HTML}{2E7D32}
\definecolor{warn-orange}{HTML}{EF6C00}
\definecolor{note-blue}{HTML}{1976D2}
\definecolor{link-blue}{HTML}{0D47A1}
\definecolor{caution-red}{HTML}{C62828}

% === GRAFICA E TABELLE ===
\usepackage{graphicx}
\usepackage{float}
\usepackage{booktabs}
\usepackage{longtable}
\usepackage{tabularx}
\usepackage{multirow}
\usepackage{array}
\usepackage{colortbl}

% === CODICE SORGENTE (Apress: grigio chiaro, numeri riga, formale) ===
\usepackage{listings}
\lstset{
  basicstyle=\footnotesize\ttfamily\color{apress-dark},
  backgroundcolor=\color{code-bg},
  frame=tb,
  framerule=0.4pt,
  rulecolor=\color{code-border},
  breaklines=true,
  breakatwhitespace=true,
  showstringspaces=false,
  tabsize=4,
  xleftmargin=2.5em,
  xrightmargin=0.5em,
  aboveskip=1em,
  belowskip=1em,
  numbers=left,
  numberstyle=\scriptsize\color{apress-lightgray},
  numbersep=10pt,
  keepspaces=true,
  keywordstyle=\color{apress-dark}\bfseries,
  stringstyle=\color{tip-green},
  commentstyle=\color{apress-lightgray}\itshape,
  columns=fullflexible,
  literate={€}{{\euro}}1 {→}{{$\rightarrow$}}1 {←}{{$\leftarrow$}}1
}
\lstdefinelanguage{yaml}{
  keywords={true,false,null,yes,no},
  sensitive=false,
  comment=[l]{\#},
  morestring=[b]",
  morestring=[b]',
}

% === INTESTAZIONI (Apress: formale con linea) ===
\usepackage{fancyhdr}
\pagestyle{fancy}
\fancyhf{}
\fancyhead[LE]{\small\itshape\leftmark}
\fancyhead[RO]{\small\itshape\rightmark}
\fancyfoot[C]{\small\thepage}
\renewcommand{\headrulewidth}{0.5pt}

% === TITOLI CAPITOLO (Apress: linea nera spessa, numero e titolo) ===
\usepackage{titlesec}
\titleformat{\chapter}[display]
  {\normalfont\bfseries}
  {\filright\sffamily\fontsize{70}{70}\selectfont\color{apress-lightgray!50}\thechapter}
  {-10pt}
  {\filright\Huge\sffamily\color{apress-dark}}
  [\vspace{2pt}{\titlerule[2pt]}]
\titlespacing*{\chapter}{0pt}{-20pt}{30pt}
\titleformat{\section}
  {\normalfont\Large\sffamily\bfseries\color{apress-dark}}
  {\thesection}{0.7em}{}
  [\vspace{1pt}{\color{apress-lightgray}\titlerule[0.3pt]}]
\titleformat{\subsection}
  {\normalfont\large\sffamily\bfseries\color{apress-gray}}
  {\thesubsection}{0.7em}{}
\titleformat{\subsubsection}
  {\normalfont\normalsize\sffamily\bfseries\color{apress-lightgray}}
  {\thesubsubsection}{0.7em}{}

% === TITOLI PARTI ===
\titleformat{\part}[display]
  {\centering\normalfont\sffamily\bfseries}
  {\fontsize{48}{48}\selectfont\color{apress-lightgray!40}\partname\ \thepart}
  {10pt}
  {\Huge\color{apress-dark}}
  [\vspace{4pt}{\titlerule[2pt]}]

% === HYPERLINK ===
\usepackage{xurl}
\usepackage[
  colorlinks=true,
  linkcolor=apress-dark,
  citecolor=apress-gray,
  urlcolor=link-blue,
  bookmarks=true,
  bookmarksnumbered=true
]{hyperref}

% === PROTEZIONE OVERFLOW ===
\tolerance=1000
\emergencystretch=1.5em
\setlength{\tabcolsep}{4pt}

% === SPAZIATURA LISTE ===
\usepackage{enumitem}
\setlist{nosep, leftmargin=2em, topsep=0.4em, itemsep=0.2em}

% === BOX ADMONITION ===
% ============================================================
% apress-boxes.tex — Box admonition stile Apress/Springer
% ============================================================
% Stile: formale, bordo nero/grigio, titoli maiuscoli,
% look accademico. Apress usa Note, Tip, Caution.
% ============================================================

\providecolor{tip-green}{HTML}{2E7D32}
\providecolor{warn-orange}{HTML}{EF6C00}
\providecolor{note-blue}{HTML}{1976D2}
\providecolor{caution-red}{HTML}{C62828}
\providecolor{apress-dark}{HTML}{1A1A1A}
\providecolor{apress-gray}{HTML}{4A4A4A}

\usepackage{fontawesome5}
\usepackage{tcolorbox}
\tcbuselibrary{skins,breakable}

% Note
\newtcolorbox{notebox}[1][]{
  colback=white, colframe=apress-dark,
  fonttitle=\sffamily\bfseries\small,
  title={\faInfoCircle\space\MakeUppercase{\LabelNote}},
  boxrule=1pt, arc=0pt, breakable,
  left=8pt, right=8pt,
  coltitle=apress-dark,
  toptitle=4pt, bottomtitle=4pt,
  separator sign none, #1
}

% Tip
\newtcolorbox{tipbox}[1][]{
  colback=tip-green!3, colframe=tip-green!60!black,
  fonttitle=\sffamily\bfseries\small,
  title={\faLightbulb\space\MakeUppercase{\LabelTip}},
  boxrule=1pt, arc=0pt, breakable,
  left=8pt, right=8pt,
  coltitle=tip-green!60!black,
  toptitle=4pt, bottomtitle=4pt,
  separator sign none, #1
}

% Warning (Caution in Apress)
\newtcolorbox{warnbox}[1][]{
  colback=warn-orange!3, colframe=warn-orange,
  fonttitle=\sffamily\bfseries\small,
  title={\faExclamationTriangle\space\MakeUppercase{\LabelCaution}},
  boxrule=1pt, arc=0pt, breakable,
  left=8pt, right=8pt,
  coltitle=warn-orange,
  toptitle=4pt, bottomtitle=4pt,
  separator sign none, #1
}

% Danger
\newtcolorbox{dangerbox}[1][]{
  colback=caution-red!3, colframe=caution-red,
  fonttitle=\sffamily\bfseries\small,
  title={\faTimesCircle\space\MakeUppercase{\LabelWarning}},
  boxrule=1pt, arc=0pt, breakable,
  left=8pt, right=8pt,
  coltitle=caution-red,
  toptitle=4pt, bottomtitle=4pt,
  separator sign none, #1
}

% Checkpoint
\newtcolorbox{checkpointbox}[1][]{
  colback=white, colframe=apress-gray,
  fonttitle=\sffamily\bfseries\small,
  title={\faCheckSquare\space\MakeUppercase{\LabelCheckpoint}},
  boxrule=1pt, arc=0pt, breakable,
  left=8pt, right=8pt,
  coltitle=apress-gray,
  toptitle=4pt, bottomtitle=4pt,
  separator sign none, #1
}

% Version
\newtcolorbox{versionbox}[1][]{
  colback=white, colframe=apress-gray,
  fonttitle=\sffamily\bfseries\small,
  title={\faCogs\space\MakeUppercase{\LabelVersionNote}},
  boxrule=0.5pt, arc=0pt, breakable,
  left=8pt, right=8pt,
  coltitle=apress-gray,
  toptitle=4pt, bottomtitle=4pt,
  separator sign none, #1
}

% Cost
\newtcolorbox{costbox}[1][]{
  colback=white, colframe=apress-gray,
  fonttitle=\sffamily\bfseries\small,
  title={\faCoins\space\MakeUppercase{\LabelCost}},
  boxrule=0.5pt, arc=0pt, breakable,
  left=8pt, right=8pt,
  coltitle=apress-gray,
  toptitle=4pt, bottomtitle=4pt,
  separator sign none, #1
}

% Exercise
\newtcolorbox{praticabox}[1][]{
  colback=white, colframe=apress-dark,
  fonttitle=\sffamily\bfseries\small,
  title={\faPencilAlt\space\MakeUppercase{\LabelExercise}},
  boxrule=1pt, arc=0pt, breakable,
  left=8pt, right=8pt,
  coltitle=apress-dark,
  toptitle=4pt, bottomtitle=4pt,
  separator sign none, #1
}


% === EURO ===
\newcommand{\euro}{€}

% === INDICE ANALITICO ===
\usepackage{makeidx}
\makeindex
